%% ma: L10-B13-0007
%% lop: 10
%% chuong: 5
%% bai: 13
%% dang: 10.13.3
%% loai: TL
%% muc: NB
%% dap_an: "-"
%% nguon: "chuyen-de-cac-so-dac-trung-cua-mau-so-lieu-khong-ghep-nhom-toan-10-chuong-trinh-moi.pdf (GV. Phan Nhật Linh) (P02), Bài 02 Dạng 1 Bài tập tự luận 7"
%% kien_thuc: "Lớp 10 Bài 13 (trung vị)"
%% trang_thai: chinh-thuc
%% ghi_chu: "Tự duyệt theo yêu cầu giáo viên 22/10/2026."
\begin{ex}
Cân nặng (kg) của một nhóm học sinh lớp 10 được cho như sau:
\begin{center}
\begin{tabular}{|c|c|c|c|c|c|c|c|c|c|}\hline
$44$ & $46$ & $46$ & $48$ & $52$ & $60$ & $75$ & $76$ & $76$ & $78$\\\hline
\end{tabular}
\end{center}
Tìm trung vị của mẫu số liệu.
\loigiai{Dãy đã sắp xếp không giảm, có $10$ giá trị nên $M_e=\dfrac{52+60}2=56$ (kg).}
\end{ex}
